\documentclass[10pt]{beamer}
\usetheme{metropolis}
\usepackage{graphicx}
\usepackage{listings}
\usepackage{booktabs}
\usepackage{xcolor}
\definecolor{codebg}{HTML}{F5F5F5}
\definecolor{codegreen}{HTML}{2E7D32}
\definecolor{codegray}{HTML}{757575}
\definecolor{codeblue}{HTML}{1565C0}
\lstdefinestyle{rstyle}{language=R,backgroundcolor=\color{codebg},basicstyle=\ttfamily\scriptsize,
  keywordstyle=\color{codeblue}\bfseries,stringstyle=\color{codegreen},
  commentstyle=\color{codegray}\itshape,breaklines=true,frame=single,
  rulecolor=\color{codegray!30},showstringspaces=false,tabsize=2,
  morekeywords={library,ggplot,aes,GGally,ggparcoord,ggradar,
    scale_color_viridis,theme_minimal,labs,coord_radar,geom_polygon,
    geom_point,geom_line,rescale}}
\lstdefinestyle{pythonstyle}{language=Python,backgroundcolor=\color{codebg},basicstyle=\ttfamily\scriptsize,
  keywordstyle=\color{codeblue}\bfseries,stringstyle=\color{codegreen},
  commentstyle=\color{codegray}\itshape,breaklines=true,frame=single,
  rulecolor=\color{codegray!30},showstringspaces=false,tabsize=4,
  morekeywords={import,from,as,pd,np,plt,parallel_coordinates,
    MinMaxScaler,polar,fill,plot}}
\title{Parallel Coordinates \& Radar Charts}
\subtitle{Workshop 5 --- Module 2}
\author{Prof.Asoc.\ Endri Raco}
\institute{Data Visualization for Data Scientists \\ UNYT Tirana}
\date{2025/26}
\begin{document}

\begin{frame}
  \titlepage
\end{frame}

\begin{frame}{Module Roadmap}
  \begin{enumerate}
    \item Parallel coordinates: each axis = one variable
    \item Reading patterns: clusters, crossings, and outliers
    \item Best practices: axis order, colour, highlighting
    \item Radar/spider charts: when they work ($\leq$3 profiles)
    \item When radar fails: area distortion and spaghetti
    \item Choosing: parallel coordinates vs radar vs pairs plot
  \end{enumerate}
  \begin{exampleblock}{Learning Outcome}
    You will produce parallel coordinate plots for 5--20 variables with
    colour-grouped highlighting, radar charts for profile comparisons
    ($\leq$3 entities), and correctly choose between the two based on
    the number of profiles and variables.
  \end{exampleblock}
\end{frame}

\section{Parallel Coordinates}

\begin{frame}{Each Axis = One Variable}
  \begin{center}
    \includegraphics[width=0.88\textwidth]{figures_w05m02/parcoord}
  \end{center}
\end{frame}

\begin{frame}{How to Read Parallel Coordinates}
  \centering
  \begin{tabular}{lp{6.5cm}}
    \toprule
    \textbf{Pattern} & \textbf{Meaning} \\
    \midrule
    Parallel lines (same slope between axes) & Positive correlation between those variables \\
    Crossing lines (X pattern) & Negative correlation --- values flip ranks \\
    Cluster of lines (same colour, similar path) & A group with similar multivariate profile \\
    Single line far from the cluster & Multivariate outlier \\
    Bottleneck (all lines pass through one point) & That variable has low variance (near-constant) \\
    \bottomrule
  \end{tabular}
  \vspace{6pt}
  \begin{alertblock}{Pro-Tip}
    Adjacent axes are compared most easily (Gestalt proximity). Place
    the two variables you most want to compare \textbf{next to each
    other}. Axis reordering is the single most impactful design decision
    in parallel coordinates.
  \end{alertblock}
\end{frame}

\begin{frame}{Best Practices: Grey + Accent}
  \begin{center}
    \includegraphics[width=0.92\textwidth]{figures_w05m02/parcoord_practices}
  \end{center}
\end{frame}

\begin{frame}[fragile]{Parallel Coordinates in Code}
  \begin{columns}[T]
    \begin{column}{0.48\textwidth}
      \textbf{R (GGally)}
\begin{lstlisting}[style=rstyle]
library(GGally)

ggparcoord(ecar,
  columns = c("FICO", "Rate",
    "Amount", "Term", "Spread"),
  groupColumn = "Tier",
  scale = "uniminmax",  # 0-1
  alphaLines = 0.15,
  showPoints = FALSE,
  order = "anyClass") + # reorder
  scale_color_viridis_d() +
  theme_minimal() +
  labs(title = "Parallel Coords",
       color = "Tier")

# Highlight one group
ggparcoord(ecar,
  columns = 1:5,
  groupColumn = "highlight",
  alphaLines = 0.05) +
  scale_color_manual(values =
    c("TRUE"="#E53935",
      "FALSE"="#CCCCCC"))
\end{lstlisting}
    \end{column}
    \begin{column}{0.48\textwidth}
      \textbf{Python (manual or pandas)}
\begin{lstlisting}[style=pythonstyle]
from sklearn.preprocessing import (
    MinMaxScaler)

cols = ["FICO","Rate","Amount",
        "Term","Spread"]
scaler = MinMaxScaler()
normed = scaler.fit_transform(
    df[cols])

fig, ax = plt.subplots(figsize=(10,5))
for i in range(min(200, len(df))):
    ax.plot(range(len(cols)),
        normed[i],
        color=cmap(norm(df["Tier"][i])),
        alpha=0.12, lw=0.6)

ax.set_xticks(range(len(cols)))
ax.set_xticklabels(cols)

# pandas built-in (simpler)
from pandas.plotting import (
    parallel_coordinates)
parallel_coordinates(df_normed,
    class_column="Tier",
    colormap="viridis", alpha=0.2)
\end{lstlisting}
    \end{column}
  \end{columns}
\end{frame}

\begin{frame}{Axis Reordering: The Key Design Decision}
  \centering
  \begin{tabular}{lp{5.5cm}}
    \toprule
    \textbf{Strategy} & \textbf{When to Use} \\
    \midrule
    Domain order & When variables have a natural sequence (e.g., time steps, pipeline stages) \\
    Correlation order & Place highly correlated pairs adjacent to reveal parallel/crossing patterns \\
    Variance order & Place highest-variance axes at edges for visual anchoring \\
    Algorithmic (clusterpath) & \texttt{ggparcoord(order="anyClass")} finds the order that best separates groups \\
    \bottomrule
  \end{tabular}
  \vspace{6pt}
  {\small Try at least 2--3 orderings before settling. The ``right''
  order depends on the question you are investigating.}
\end{frame}

\section{Radar Charts}

\begin{frame}{Radar: Profile Comparison ($\leq$3 Profiles)}
  \begin{center}
    \includegraphics[width=0.55\textwidth]{figures_w05m02/radar}
  \end{center}
\end{frame}

\begin{frame}{When Radar Fails}
  \begin{center}
    \includegraphics[width=0.88\textwidth]{figures_w05m02/radar_fails}
  \end{center}
  \begin{alertblock}{Pro-Tip}
    Radar charts have three known problems: (1) \textbf{area is
    misleading} --- the enclosed area depends on axis order, not just
    the values. (2) \textbf{Axis order is arbitrary} for non-sequential
    variables, so the shape is meaningless. (3) Beyond 3 profiles,
    they become spaghetti. Use radar only for \textbf{2--3 profiles
    with $\leq$8 axes} where the goal is holistic comparison, not
    precise value reading.
  \end{alertblock}
\end{frame}

\begin{frame}[fragile]{Radar in Code}
  \begin{columns}[T]
    \begin{column}{0.48\textwidth}
      \textbf{R (ggradar or manual)}
\begin{lstlisting}[style=rstyle]
# ggradar (simple)
# install.packages("ggradar")
library(ggradar)
ggradar(profiles,
  grid.min = 0, grid.max = 1,
  group.line.width = 1.5,
  group.point.size = 3,
  fill = TRUE, fill.alpha = 0.1,
  legend.position = "bottom")

# Manual with coord_polar
ggplot(df_long, aes(x = variable,
    y = value, group = profile,
    color = profile)) +
  geom_polygon(fill = NA,
    linewidth = 1.2) +
  geom_point(size = 2) +
  coord_polar() +
  scale_color_manual(values =
    c("#1565C0","#E53935")) +
  theme_minimal()
\end{lstlisting}
    \end{column}
    \begin{column}{0.48\textwidth}
      \textbf{Python (manual polar)}
\begin{lstlisting}[style=pythonstyle]
categories = ["FICO","Rate",
    "Amount","Term","Spread","Appr"]
N = len(categories)
angles = [n/N * 2 * np.pi
          for n in range(N)]
angles += angles[:1]  # close polygon

fig, ax = plt.subplots(
    subplot_kw=dict(polar=True))

for profile, color, label in [
    (vals1, "#1565C0", "Tier 1"),
    (vals2, "#E53935", "Tier 5")]:
    values = profile + profile[:1]
    ax.fill(angles, values,
        alpha=0.1, color=color)
    ax.plot(angles, values, "-o",
        color=color, lw=2,
        label=label)

ax.set_xticks(angles[:-1])
ax.set_xticklabels(categories)
ax.legend()
\end{lstlisting}
    \end{column}
  \end{columns}
\end{frame}

\section{Choosing the Right Chart}

\begin{frame}{Parallel Coordinates vs Radar vs Pairs}
  \begin{center}
    \includegraphics[width=0.92\textwidth]{figures_w05m02/comparison}
  \end{center}
\end{frame}

\begin{frame}{Decision Rules}
  \centering
  \begin{tabular}{lll}
    \toprule
    \textbf{Question} & \textbf{Answer} & \textbf{Chart} \\
    \midrule
    How many profiles to compare? & 2--3 & Radar \\
    How many profiles? & Many (100s+) & Parallel coordinates \\
    How many variables? & 3--8 & Radar or pairs \\
    How many variables? & 8--20 & Parallel coordinates \\
    Need precise value reading? & Yes & Pairs plot (scatter) \\
    Need holistic ``shape'' comparison? & Yes & Radar \\
    Need group separation? & Yes & Parallel coords + colour \\
    \bottomrule
  \end{tabular}
  \vspace{6pt}
  \begin{exampleblock}{Data Insight}
    In practice, \textbf{parallel coordinates} are far more common in
    data science than radar charts. Radar charts are popular in sports
    analytics (player profiles) and product comparisons but are
    statistically problematic. Default to parallel coordinates unless
    you have exactly 2--3 profiles for a general audience.
  \end{exampleblock}
\end{frame}

\section{Interactive Extensions}

\begin{frame}[fragile]{Interactive Parallel Coordinates}
\begin{lstlisting}[style=rstyle]
# R: plotly (brushable axes)
library(plotly)
plot_ly(ecar, type = "parcoords",
  line = list(color = ~Tier,
    colorscale = "Viridis"),
  dimensions = list(
    list(label = "FICO", values = ~FICO),
    list(label = "Rate", values = ~Rate),
    list(label = "Amount", values = ~Amount),
    list(label = "Term", values = ~Term),
    list(label = "Spread", values = ~Spread)))
\end{lstlisting}
\begin{lstlisting}[style=pythonstyle]
# Python: plotly
import plotly.express as px
fig = px.parallel_coordinates(
    ecar, color="Tier",
    dimensions=["FICO","Rate","Amount","Term","Spread"],
    color_continuous_scale="Viridis")
fig.show()  # interactive: brush each axis to filter
\end{lstlisting}
  \begin{alertblock}{Pro-Tip}
    Interactive parallel coordinates (plotly) allow \textbf{brushing}:
    drag on any axis to select a range, and all lines outside that
    range fade. This is the most powerful exploratory use of parallel
    coordinates --- impossible in static plots.
  \end{alertblock}
\end{frame}

\section{Synthesis}

\begin{frame}{Key Takeaways}
  \begin{enumerate}
    \item \textbf{Parallel coordinates}: each vertical axis = one
          variable; lines connect one observation across all axes.
          Best for 5--20 variables, many observations.
    \item \textbf{Reading patterns}: parallel = positive correlation,
          crossing = negative, cluster = group, stray line = outlier.
    \item \textbf{Axis order is critical}: place correlated pairs
          adjacent. Try \texttt{ggparcoord(order="anyClass")} for
          algorithmic reordering.
    \item \textbf{Grey+accent} prevents spaghetti: draw all lines grey,
          highlight one group in colour.
    \item \textbf{Radar charts} work only for $\leq$3 profiles,
          $\leq$8 axes. Area is misleading; axis order is arbitrary.
    \item \textbf{Interactive} parallel coordinates (plotly) with
          axis brushing is the most powerful multivariate exploration
          tool available.
  \end{enumerate}
\end{frame}

\begin{frame}{Essential Readings}
  \begin{itemize}
    \item Inselberg, A. (2009). \emph{Parallel Coordinates: Visual
          Multidimensional Geometry}. Springer.
    \item Wilke, C.~O. (2019). \emph{Fundamentals of Data Visualization},
          Chapter 6 (Parallel Coordinates).
    \item Few, S. (2006). ``Radar Graphs: Not So Useful After All.''
          Visual Business Intelligence Newsletter.
    \item plotly documentation: \texttt{https://plotly.com/r/parallel-coordinates-plot/}
  \end{itemize}
  \vspace{6pt}
  \textbf{Next Module}: Dimensionality Reduction Visualization ---
  PCA, t-SNE, UMAP (W05-M03)
\end{frame}

\end{document}
